\subsection{Retaining load dependencies after process discovery}
All 288 fits complete. Every routing-feedback calibration and validation episode satisfies the pre-crisis criterion. One of the 3168 episodes, in P4 combined feedback, violates it; that dataset remains in the diagnostic results.

With automatically selected structure, retained-load planning reduces the equal-process mean normalized cost from 0.627 to 0.542, a 13.5\% reduction relative to predictive selection. The primary paired difference is $-0.0848$, with multiplicity-adjusted 97.5\% interval $[-0.1353,-0.0343]$. Table~\ref{tab:structures_results} shows the individual structures and stronger comparators. On their original cost scales, mean cost decreases in P1, P2, P3, P5 and P6 by 11.6\%, 14.5\%, 7.0\%, 31.0\% and 24.9\%, respectively. Among the six secondary process contrasts, P6 retains a nonzero benefit after correction, with difference interval $[-18.48,-1.82]$ thousand. P4 instead increases by 16.8\%, with adjusted interval $[2.30,8.74]$ thousand. 

\begin{table}[t]\caption{Routing feedback, selected discovered structure: mean per-fit cost divided by matched no-action cost. Regret is retained-load minus known-generator cost, normalized by no action. Each row contains 12 independent fits.}\label{tab:structures_results}\centering\small\setlength{\tabcolsep}{4pt}
\begin{tabular}{lrrrrrr}\toprule
Process & Static & Predictive & Retained & Heuristic & Generator & Regret\\\midrule
P1 & 0.584 & 0.555 & 0.491 & 0.499 & 0.429 & 0.062 \\
P2 & 0.688 & 0.658 & 0.561 & 0.648 & 0.435 & 0.126 \\
P3 & 0.659 & 0.659 & 0.613 & 0.635 & 0.522 & 0.092 \\
P4 & 0.518 & 0.527 & 0.615 & 0.623 & 0.513 & 0.102 \\
P5 & 1.030 & 0.828 & 0.571 & 0.467 & 0.415 & 0.155 \\
P6 & 0.573 & 0.537 & 0.404 & 0.574 & 0.351 & 0.054 \\
\bottomrule\end{tabular}\end{table}


The control mechanisms identify when component retention helps. The descriptive retained-minus-predictive normalized difference is $-0.055$ under combined feedback, but $+0.050$ under fixed behavior and $+0.110$ under service feedback. Retaining weakly estimated return dependencies is useful in the routing comparison, but can add error when those dependencies are absent. Figure~\ref{fig:structural_matrix} separates these mechanisms and the two discovery procedures.

\begin{figure}[t]\centering
\includegraphics[width=\textwidth]{generated/structural_matrix_en.pdf}
\caption{Mean cost reduction of retained-load relative to predictive selection within each structural variant, in percent. Positive values favor retention. Both discovery procedures use the same corrected executor. Cell estimates are descriptive; the primary aggregate comparisons are reported in the text.}\label{fig:structural_matrix}
\end{figure}

\subsection{Executable structure preserves intervention choices}
The selected discovered models have the reference visible language in all 288 fits, with mean training precision between 0.996 and 1.000. They reproduce the known-structure planning costs exactly for all four duration/routing component combinations. Standard IM recovers the reference language in 144 fits; P1, P2 and P6 admit additional paths. In the second primary comparison, changing from standard IM to selected discovery reduces retained-load normalized cost from 0.638 to 0.542, or 15.0\%. The paired difference is $-0.0958$, with adjusted 97.5\% interval $[-0.1205,-0.0711]$.

P2 illustrates the consequence for response choice. Selected discovery chooses rework reduction in 8 of 12 fits, whereas standard IM chooses capacity in every fit; normalized costs are 0.561 and 0.898. Both variants use identical learned behavior and the same corrected executor. Preserving return structure and observed route constraints changes which action that behavior supports. Reproducing the training traces alone is therefore insufficient to validate a model for intervention planning.

\subsection{System-wide planning and simple feedback rules}
P6 routing feedback supplies a diagnostic example of simulation improving on the local feedback rule: normalized costs are 0.404 and 0.574, with descriptive paired 95\% difference interval $[-24.85,-0.47]$ thousand. The rule chooses branch-specific capacity in 4 fits; retained-load chooses rework or the shared pool. The comparisons for P3 and P5 do not establish an advantage over the rule. Under P6 service feedback, retained-load also beats the rule in mean cost, 0.671 versus 0.878, but static planning costs only 0.496. Simulation can help assess system-wide consequences without requiring every load component.

The resource replay shows why those consequences matter. Under P5 combined feedback, adding a B worker reduces its mean queue from 137.95 to 37.29 while increasing D's from 1.28 to 116.13; the dominant queue moves in all 24 matched diagnostic paths. Exact control flow does not, however, guarantee accurate behavioral extrapolation. In P4 routing feedback, removing the learned return-load terms reduces normalized cost from 0.615 to 0.514; removing the duration-load terms leaves it at 0.615. This ablation locates the decision error in the estimated return dependencies.

\subsection{Sensitivity to selection and intervention assumptions}
In P1, P2, P3, P5 and P6, retained-load has lower mean routing cost at every tested setting of either selector threshold. Varying $\Delta_R$ from 0 to 0.01 changes P5's predictive-selected normalized cost from 0.726 to 0.928, versus 0.571 for retained-load. P4 favors predictive selection throughout. The effect extends beyond the default selector cutoff, while remaining process-dependent. These sweeps diagnose the fixed methods; they do not choose a new threshold from test outcomes.

With P2's complex-case fraction shifted to 65\%, the unchanged retained and predictive policies cost 0.668 and 0.679. In P3, actual rework reductions of 30\% yield costs 0.730 and 0.683, whereas reductions of 70\% yield 0.518 and 0.639. The ordering reverses without reselection. A prepared response therefore needs an explicit range of intervention effects under which its recommendation remains useful.
